\documentclass[12pt]{article}
\usepackage[utf8]{inputenc}
\usepackage[T1]{fontenc}
\usepackage{hyperref}
\title{Guibian Zhuang Language Resources for K-12 Curriculum Development}
\author{Compiled by Claude Sonnet 5.5 (claude-sonnet-5-5) for Andr\'as Kornai}
\date{}
\begin{document}
\maketitle

\section{Introduction}
Guibian Zhuang (ISO 639-3 zgn; Glottocode guib1244) is a Northern Zhuang variety in the Tai-Kadai family, spoken in China. Glottolog lists it as not endangered, at EGIDS level 3 (wider communication). Speaker figures differ widely between sources: Global Recordings Network gives 100,000 and Joshua Project about 1,091,000, while the task description gives about 420,000. Alternate names include Buyei, Kang Yei and Northern Zhuang, so care is needed not to confuse it with Buyei in Guizhou, which is treated as a separate language in China. No school materials specific to this variety were found in the searches.

\section{Primary Educational Resources}
No primers, textbooks or school programmes specific to Guibian Zhuang were found. Standard Zhuang (code zha) materials may be relevant, since it is based on Northern Zhuang.

\section{Linguistic and Reference Materials}
\begin{itemize}
\item Glottolog entry for Guibian Zhuang. \url{https://glottolog.org/resource/languoid/id/guib1244}
\item Glottolog entry by ISO code. \url{https://glottolog.org/resource/languoid/iso/zgn}
\item Wikitongues language page. \url{https://wikitongues.org/languages/zgn/}
\item Detect Language, Guibian Zhuang overview. \url{https://detectlanguage.com/languages/zgn}
\end{itemize}

\section{Religious and Community Resources}
\begin{itemize}
\item Joshua Project, Guibian Zhuang in China (people group profile). \url{https://joshuaproject.net/people_groups/19808}
\end{itemize}

\section{Digital and Audio Resources}
\begin{itemize}
\item Global Recordings Network, ``Zhuang, Guibian language'' (audio recordings listing). \url{https://globalrecordings.net/en/language/zgn}
\end{itemize}

\section{Curriculum Development Context}
The searches found no Guibian-specific curriculum resources. Any programme would need to decide whether to build on Standard Zhuang materials or on local varieties, and to settle the relationship with neighbouring Buyei. Local community involvement is essential.

\end{document}